\documentclass[11pt,a4paper]{article}
\usepackage[T1]{fontenc}
\usepackage[margin=2.7cm]{geometry}
\usepackage{amsmath,amssymb,amsthm}
\usepackage{booktabs}
\usepackage[hidelinks]{hyperref}

\newtheorem{lemma}{Lemma}
\newcommand{\dt}{\Delta t}
\newcommand{\R}{\mathbb{R}}

\title{Order reduction of Gau{\ss} collocation for the HDG Timoshenko wave solver\\
  \large temporal order $s{+}1$ under time-dependent Dirichlet data,
  full order $2s$ under compatible data}
\author{ne9 notes}
\date{July 16, 2026}

\begin{document}
\maketitle

\section{Summary of observations}

The condensed HDG Timoshenko wave solver is stepped by $s$-stage Gau{\ss}
collocation (temporal design order $2s$; only $\lceil s/2\rceil$ complex
eigen-representatives are solved per step). Temporal convergence sweeps with the
spatial discretization held constant (degree $5$, $n_x = 256$, cross domain,
experiment \texttt{ne9-11-conv-t-stages.sh}) give the orders in
Table~\ref{tab:orders}. The two arms use the \emph{same} manufactured solution
family (\texttt{wave4}), differing only in a spatial phase $p_x$ of the
standing-wave factor $\cos(\omega(x{+}y{+}z) + p_x)$:
\begin{itemize}
  \item \textbf{stiff arm} ($p_x = 0$): the Dirichlet data at the clamped tips is
    $g(t) \propto \cos(\omega t + p)$ --- the supports move;
  \item \textbf{compatible arm} ($p_x = -\pi/2$): the spatial factor becomes
    $\sin(\omega s)$, which vanishes at every tip, so $g \equiv 0$ for all $t$
    while the interior solution (both $u$ and $r$, all time phases) is unchanged
    in structure.
\end{itemize}

\begin{table}[h]\centering
\begin{tabular}{c c cccc c cccc}
\toprule
 & & \multicolumn{4}{c}{stiff arm ($p_x = 0$)} & \phantom{x}
   & \multicolumn{4}{c}{compatible arm ($p_x = -\pi/2$)} \\
\cmidrule{3-6}\cmidrule{8-11}
$s$ & $n_t$: & 8 & 16 & 32 & 64 & & 8 & 16 & 32 & 64 \\
\midrule
1 & & 1.84 & 1.96 & 2.01 & 2.00 & & 1.47 & 1.88 & 1.98 & 1.99 \\
2 & & 3.31 & 3.14 & 3.03 & 2.50 & & 3.83 & 3.96 & 4.00 & 4.00 \\
3 & & 4.02 & 4.00 & 4.01 & 4.00 & & 5.89 & 5.97 & 6.00 & 6.00 \\
\bottomrule
\end{tabular}
\caption{Observed temporal orders $\log_2(e_{n_t/2}/e_{n_t})$ of the relative
  $L^2$ displacement error, at fixed spatial discretization (deg $5$,
  $n_x = 256$). Compatible data reaches the design order $2s$; moving Dirichlet
  data caps the scheme at its \emph{stiff} order: $s{+}1$ for odd $s$, degrading
  towards $s$ for even $s$.}
\label{tab:orders}
\end{table}

This is \emph{not} an implementation defect: the tableau, the eigenbasis stage
protocol and the value-form endpoint recombination reproduce order $2s$ exactly
in a nonstiff scalar replication, and the identical temporal errors are obtained
for different spatial degrees and resolutions
(\texttt{experiments/gauss\_value\_form\_check.py}). It is the classical
\emph{order reduction} of implicit Runge--Kutta methods on stiff problems
\cite{prothero,hw2,ssvh}.

\section{Mechanism: every eigenmode is a Prothero--Robinson problem}

After eliminating the algebraic variables, the semidiscrete system is
$\dot x = L_h x + B_h\, g(t) + (\text{interior loads})$. The spectrum of the
first-order wave operator reaches $|\lambda_{\max}| \sim c\,(p{+}1)/h$, so
$\dt\,|\lambda_{\max}|$ is (up to constants) the CFL number: implicit time
stepping above the CFL limit \emph{is} the stiff regime. Expanding
$x = \sum_k y_k v_k$ in the eigenbasis $L_h v_k = \lambda_k v_k$
($\lambda_k = i\omega_k$) gives scalar equations, all driven by the same data,
\begin{equation*}
  \dot y_k = \lambda_k y_k + b_k\, g(t),
  \qquad b_k = \text{(modal coefficient of the boundary input)} .
\end{equation*}
Let $\varphi_k$ be the forced response, $\dot\varphi_k = \lambda_k\varphi_k +
b_k g$. Substituting this identity puts the mode in Prothero--Robinson normal
form \cite{prothero}
\begin{equation}
  \dot y_k = \lambda_k\,(y_k - \varphi_k) + \dot\varphi_k ,
  \qquad y_k \equiv \varphi_k \ \text{if started on the forced trajectory.}
  \label{eq:pr}
\end{equation}
For stiff modes the forced response is the quasi-static slaving series
\begin{equation}
  \varphi_k(t) = -\frac{b_k}{\lambda_k}\,g(t)
                 -\frac{b_k}{\lambda_k^2}\,g'(t) - \cdots ,
  \qquad \partial_t^{\,j}\varphi_k \propto g^{(j)} ,
  \label{eq:slaving}
\end{equation}
and smoothness of the exact solution forces its stiff-mode content to be exactly
\eqref{eq:slaving} (any free oscillation $e^{i\omega_k t}$ would destroy
smoothness of time derivatives). Collocation applied to \eqref{eq:pr} in the
stiff limit $|\lambda_k\dt| \gg 1$ has the endpoint error recursion
\cite{hw2}
\begin{equation}
  e_{n+1} = R(\lambda_k\dt)\, e_n + \delta_n, \qquad
  \delta_n = C\,\dt^{\,s+1}\varphi_k^{(s+1)}(\xi_n) + O(\dt^{\,s+2}),
  \label{eq:defect}
\end{equation}
i.e.\ the leading defect is proportional to the first time derivative of the
tracked data that the degree-$s$ collocation polynomial cannot represent. For
Gau{\ss} methods $|R(iy)| = 1$ (unitary) and $R(\infty) = (-1)^s$: nothing damps
the defects; they accumulate with alternating sign for odd $s$ (telescoping to
global order $s{+}1$) and with constant sign for even $s$ (order $s$). This
even/odd asymmetry matches Table~\ref{tab:orders}, including the degradation of
$s = 2$ towards order $2$ as the sweep descends deeper into the stiff regime.

\section{Why boundary data excites the stiff modes (and interior loads do not)}

The reduction requires $O(1)$ total weight $b_k$ across the stiff spectrum.
Homogenizing $u(0,t) = g(t)$ with a lifting, e.g.\ $v = u - (1-x)\,g(t)$ in 1D,
yields a homogeneous-BC problem with source $-(1-x)\,\ddot g(t)$. The profile
$(1-x)$ is smooth but \emph{violates the boundary condition}; integration by
parts against the eigenfunctions leaves a boundary term, so its eigencoefficients
decay only like $1/k$ --- every stiff mode is driven. (In the HDG scheme the
data actually enters through the trace pinning on the boundary-adjacent elements,
an input of support $O(h)$: a boundary delta layer in the limit, with the same
conclusion.) In contrast, a spatially smooth interior load that is compatible
with the boundary conditions has spectrally decaying coefficients: the stiff
modes' forced responses \eqref{eq:slaving} are negligible and the classical
order-$2s$ theory applies. Summed over modes, the surviving boundary-driven
defect \eqref{eq:defect} carries an $O(1)$ constant; parabolic problems soften
this to fractional orders $s{+}1{+}\alpha$ through analytic smoothing
($|R| < 1$ off the imaginary axis) \cite{or}, while the unitary wave case shows
the clean integer $s{+}1$ observed here.

\section{The compatibility condition}

\begin{lemma}[per-step exactness]
If the tracked data $\varphi$ is a polynomial of degree $\le s$ on a step, the
collocation step for \eqref{eq:pr} is exact for every $\lambda$: the error
$e = u - \varphi$ is a degree-$s$ polynomial with $e(t_n) = 0$ and
$e'(t_j) = \lambda\, e(t_j)$ at the $s$ Gau{\ss} nodes, a homogeneous collocation
problem whose only solution is $e \equiv 0$ (away from a measure-zero resonance
set of $\lambda\dt$).
\end{lemma}

By \eqref{eq:slaving}, $\partial_t^{\,j}\varphi_k \propto g^{(j)}$, so the
operative compatibility condition on the boundary data is simply
\begin{equation*}
  \boxed{\ g^{(s+1)} \equiv 0 \ }
  \qquad\text{(per step; e.g.\ $g$ globally polynomial of degree $\le s$),}
\end{equation*}
which also covers the homogeneous case $g \equiv 0$ and time-independent
inhomogeneous data. Equivalently, in the language of \cite{ssvh,or}: the time
derivatives of the solution from order $s{+}1$ upwards must satisfy the
(homogeneous) boundary conditions, so that $\|L_h^m\, \partial_t^{\,j} x\|$
stays bounded uniformly in $h$. Note this is stronger than the textbook
\emph{corner} conditions at $t = 0$: the defect re-enters at every step, so the
conditions must hold for all $t$. Smoothness of $g$ does not help --- $\cos
\omega t$ is analytic and maximally incompatible.

Inhomogeneous compatible tests are easy to construct by linearity: superposing a
zero-forcing static field ($u \mathrel{+}= c\,s$ per arm, constant tip data) or
a ramp ($u \mathrel{+}= (a + bt)\,c\,s$, tip data linear in $t$) onto the
compatible arm leaves the forcing untouched and, by the Lemma, contributes no
temporal error at all.

\section{Consequences}

\begin{itemize}
\item The engineering-relevant regime --- clamped supports, interior/inertial
  loading --- is the compatible one: the solver delivers its full design order
  $2s$ there with $\lceil s/2\rceil$ global solves per step and no
  post-processing. (An endpoint re-solve of the algebraic variables was
  prototyped and removed again: for compatible data the value-form recombination
  of trace and dual is already \emph{exact} --- linear time-invariant slaving of
  a degree-$s$ polynomial --- and the extra recovery amplified round-off by
  $\sim h^{-4}$ in the trace error.)
\item Boundary-driven problems (prescribed support motion) are capped at order
  $s{+}1$ / $s$ regardless of spatial resolution --- refining in space makes the
  system stiffer, not better. If full order ever matters there, the known remedy
  is corrected stage boundary values \cite{amp,cgad}, not more stages or finer
  meshes; matching temporal to spatial order via odd stage counts ($s = p$) is
  the brute-force alternative.
\end{itemize}

\begin{thebibliography}{9}\small
\bibitem{prothero} A.~Prothero, A.~Robinson,
  \emph{On the stability and accuracy of one-step methods for solving stiff
  systems of ordinary differential equations}, Math.\ Comp.\ 28 (1974).
\bibitem{hw2} E.~Hairer, G.~Wanner,
  \emph{Solving Ordinary Differential Equations II}, Springer,
  \S IV.15 (B-convergence, Prothero--Robinson analysis).
\bibitem{ssvh} J.~M.~Sanz-Serna, J.~G.~Verwer, W.~H.~Hundsdorfer,
  \emph{Convergence and order reduction of Runge--Kutta schemes applied to
  evolutionary problems in partial differential equations},
  Numer.\ Math.\ 50 (1986).
\bibitem{or} A.~Ostermann, M.~Roche,
  \emph{Runge--Kutta methods for partial differential equations and fractional
  orders of convergence}, Math.\ Comp.\ 59 (1992).
\bibitem{cgad} M.~H.~Carpenter, D.~Gottlieb, S.~Abarbanel, W.-S.~Don,
  \emph{The theoretical accuracy of Runge--Kutta time discretizations for the
  initial boundary value problem}, SIAM J.\ Sci.\ Comput.\ 16 (1995).
\bibitem{amp} I.~Alonso-Mallo, C.~Palencia,
  \emph{Optimal orders of convergence for Runge--Kutta methods and linear,
  initial boundary value problems}, Appl.\ Numer.\ Math.\ (2003); see also
  I.~Alonso-Mallo, B.~Cano on avoiding order reduction via corrected boundary
  values.
\end{thebibliography}

\emph{Reproducibility:} \texttt{experiments/ne9-11-conv-t-stages.sh} (both arms),
rates via \texttt{experiments/ne9-conv-rates.sh <json> nt gauss\_stages};
scalar mechanism isolation: \texttt{experiments/gauss\_value\_form\_check.py}.
Bibliographic details are from memory and should be verified before reuse.

\end{document}
